\section{Experimental Design}
\label{sec:experimental-design}

The aggregate LOS experiments incorporate Virna Lissi Mart\'inez S\'anchez's
report and structural sensitivity analysis \citep{martinez2026los}. The
room-execution campaigns are evaluated separately. Tables and figures are
regenerated from pinned aggregate results; original contributions retain
their attribution and source hashes.

\subsection{Data and predictive modelling}
\label{sec:data-prediction}

\subsubsection{Source data and targets}

The original workbook audit describes 14\,649 cases, 30 variables and
11\,108 patients; 2546 patients have repeated cases, up to nine per person
\citep{group2026dataaudit}. Interventions span 2 January 2019--30 December
2022, with one admission in 2018. Records cover demographics, clinical codes,
admission/discharge and pathway timestamps. They do not reconstruct reliable
room/session assignments, staffing, bed capacity, cancellations or rejected
dates. Experimental calendars and capacities are therefore assumptions.
Raw, cleaned, prediction-eligible and matched cohorts have different denominators.

Targets are inclusive admission-to-discharge LOS and room entry-to-exit
occupancy, distinct from incision-to-closure time. The audit found no exact
duplicate rows, but 2560 LOS values disagree with the date-derived target;
1884 recorded one-day stays span a night, and 31 recorded values resemble
Julian dates rather than durations. A one-day value alone cannot identify
ambulatory care. Zero clocks are treated as missing unless midnight is
confirmed: the audit reports 172 room-entry, 172 room-exit, 552 incision and
133 pre-anaesthesia zeros. Inverted intervals are not automatically repaired
by adding 24 hours. Section~\ref{sec:data-audit} retains further aggregate
checks; no original patient or professional identifiers are reproduced.

\subsubsection{Prediction time and evaluation}

Temporal and patient-disjoint evaluation measure different generalisation
settings. Imputation, encoding and category
rules use training data only. Historical medians, MAE, tail errors and
calibration provide complementary references. Features must be available
at decision time: discharge and observed pathway durations are excluded;
planned diagnoses, procedures, anaesthesia and clinician information remain
conditional on local preoperative availability. A changed proposed date
also changes any date-derived predictor. The implemented notebook and room
baseline splits are distinguished below.

\subsubsection{LOS predictor and matching interface}

The contributor uses XGBoost \citep{chen2016} on $\log(1+\mathrm{LOS})$ and
reports test MAE 0.412 days, $R^2=0.789$ and 89.4\% coverage for a nominal
90\% interval with half-width 1.02 days \citep{martinez2026los}. Scheduling
CSVs do not independently verify these metrics. The bridge uses a fixed
half-width in the inductive confidence-prediction tradition
\citep{papadopoulos2002}.

Back-transformed predictions are floored at one day and matched by patient
key and normalised intervention date, averaging duplicate key-date predictions.
Reported match rates are 86.7\% for the small workload and 86.4\% for the
large one. Unmatched cases receive the median matched prediction in the
selected window, or three days if no match exists; this is not an independently
estimated training median. Point and interval endpoints are rounded for
planning. The HIGH/LOW flag thresholds regression output above four days
and is distinct from the contributor's separately evaluated risk classifier.

\subsection{LOS instances, benchmarks and scenarios}

The original campaign compares SA, tabu, the hybrid, GA and ACO on 45 cases
with 35 sessions over five dates and 420 cases with 210 sessions over
30 dates. Capacity is 480 minutes per speciality-day and 42 aggregate beds.
The large workload has a speciality requiring 15\,438 minutes against
14\,400 available, so unavoidable overload is at least 1038 minutes and the
weighted objective has lower bound 5190, without certifying its full optimum.

The structural campaign splits each daily session into two 240-minute
sessions, preserving total capacity and giving 420 sessions. It tests SA,
tabu, the hybrid and GA with mutation $1/n$, without ACO. Original GA uses
0.08; a separate sensitivity run tests $1/n$. Source stopping rules and
unequal computation are retained rather than interpreted as a common-budget
comparison.

For each method and seed, point-LOS and upper-endpoint schedules are frozen
and evaluated on 200 common scenarios. Ten seeds, 42--51, support the main
comparisons; paired SE uses ten seed-level upper-minus-point scenario means,
not 200 independent optimiser runs. Original noise is uniform in
$[-1.02,1.02]$ days around unrounded predictions, clipped at one day and
integer-rounded. Alternatives use symmetric bounded noise with stochastic
rounding and extended, rather than horizon-truncated, bed tracking. These
are sensitivity rules, not distributions implied by interval coverage.
Additional experiments vary bottleneck-speciality block load across
0.80, 0.93, 0.98 and 1.10 with five seeds, keeping 42 beds. Matched-only
analysis retains 363 cases and removes 57 fallback cases, changing demand.

\subsection{Room/Mesa cohorts and paired experiment matrix}
\label{sec:room-design}

The separate room study uses the cleaned workbook, not the raw audit's
14\,649-row denominator. Its saved manifest records 14\,507 source rows,
14\,434 usable episodes and 73 exclusions, all zero clock endpoints in that
file version. Duration medians are learned from 10\,800 valid 2019--2021
cases, with a ten-observation threshold per procedure and a 70-minute global
fallback. Evaluation uses 2022. This temporal split permits recurring patients, unlike the patient-disjoint
model-preprocessing partitions described below.

The first room campaign contains one seven-case fixture, three historical
days (3 January: 9 cases; 2 February: 15; 20 October: 25), and two generated
samples at each of 100 and 300 cases. The selected dates have no exclusions; the middle and largest workloads
follow valid daily case counts among such dates. Generated instances sample complete eligible 2022
rows with replacement, retaining procedure--duration associations but not
within-day dependence. Each workload uses six methods and solver seeds
0, 1 and 2. Its 144 initial plans are each replayed under static/reactive
policies and closure/no-closure scenarios: 576 executions, not 576 independent
hospital days \citep{group2026roomstudy}.

Historical and generated days open 08:00--17:00 with 15-minute turnover.
Historical workloads use two interchangeable rooms; room 1 forbids starts
10:00--12:00 in the closure scenario. The small fixture uses two rooms,
08:00--11:00 and a 09:00--10:00 closure. Initial plans do not know the future
closure. For generated workloads,
\begin{equation}
 |\mathcal R|=\max\!\left(1,
 \left\lceil\frac{\sum_i(p_i+\tau)}{(b-a)\rho}\right\rceil\right),
 \qquad\rho=1.1.
\end{equation}
The four samples therefore use 14, 15, 44 and 45 rooms. Integer rounding
changes achieved load. This is an overloaded target experiment, not a
light-load sensitivity campaign; a single closure also represents a smaller
capacity fraction as the room pool grows.

Each nontrivial request has 1000 fitness calls and a 60-second safety limit.
No deadline stops occur in the saved campaigns. The baseline performs no
search evaluations. Reactive execution may use two additional request
budgets, so policies do not receive equal total computation. Runtime includes worker startup and communication. The 100-case runs were
sequential and 300-case runs used four concurrent batches, limiting timing
comparisons. Manifests record Python 3.12.3, Mesa 3.5.1, dependency versions
and source fingerprints, including dirty-checkout provenance.

\subsection{Duration-information experiment and predictive interfaces}
\label{sec:duration-design}

A separate saved campaign repeats the same matrix in median and oracle modes,
providing 1152 executions from 288 initial plans. These median runs are a
rerun, not extra independent replications of the first campaign. Oracle
mode replaces only scheduling estimates with realised room-occupancy minutes.
Workloads, execution durations and resources are held fixed; generated room
counts are computed from the original medians before substitution. Initial
plans are recomputed across modes but shared within each mode's four execution
conditions. Future closures remain hidden. This is a perfect-information
reference under finite search and a fixed decoder, not an ML predictor or a
certified operational upper bound. Normalised objective costs change scale
between modes; comparisons therefore use completions and realised overtime.

The operating-duration model branch \citep{group2026durationmodel} validates
patient-disjoint chronological partitions, excludes calendar features by
default, separates selection from later patient-disjoint calibration, and
selects the fastest candidate within 0.2 minutes of the best validation MAE.
Test data are reserved for reporting. Its review labels older metric tables
as belonging to a previous protocol; those are not verified current-model
results. Neither the median nor oracle campaign evaluates a learned room
model within Mesa.

Both preprocessing notebooks use pre-2021 training, new-patient validation
in 2021 and new-patient testing from 2022, excluding later cases of patients
in earlier partitions. Split identifiers/dates remain audit metadata.
Target-derived fields and postoperative/discharge information are excluded
from predictors; clinical codes require confirmed preoperative availability.
These controls must survive model-to-scheduler integration.

\subsection{Verification and reproducible evidence}

The incorporated room tables are regenerated from pinned aggregate CSVs.
Checks enforce unique run keys, four replays per initial plan, case conservation,
paired mode coverage and equality of no-closure policy outcomes. An independent
review of all 576 primary and 1152 duration execution exports recomputes
occupancy and overtime and checks start windows, closures, turnover separation,
evaluation caps and shared initial schedules. All pass, with zero reported
coordinator/search failures, timeouts or deadline stops. Validation of the reviewed source snapshot ran 94 tests, with 92 passing
and two optional tests skipped. These
checks support implementation consistency; they do not validate clinical
assumptions or make solver-seed repetitions independent observations.


\subsection{Additional source-workbook audit details}
\label{sec:data-audit}

By admission year, the raw audit counts are 1 (2018), 3690 (2019),
3526 (2020), 3764 (2021) and 3668 (2022); intervention dates are in 2019--2022.
All interventions lie between admission and discharge in that audit, with
no missing date endpoints or discharge before admission. These checks retain the raw-audit denominator rather than that of the
cleaned room cohort.

\begin{table}[H]
\centering\small
\caption{Selected additional quality checks in the original workbook audit.
Timestamp comparisons exclude zero clocks and use their own valid-pair denominators.}
\label{tab:source-quality}
\begin{tabularx}{\linewidth}{@{}YY@{}}
\toprule
Check & Audit result \\
\midrule
Associated diagnosis fields 1--5 & Missingness 20.7\%, 43.2\%, 63.2\%, 77.1\%, 86.9\%. \\
Secondary procedure fields 2--4 & Missingness 7.3\%, 38.6\%, 66.0\%. \\
Anaesthesia / locoregional anaesthesia & Missingness 15.3\% / 37.1\%. \\
Room exit before entry & 0 of 14\,477 comparable pairs. \\
Room exit before incision / incision before entry & 4 / 42 of 14\,026 pairs. \\
Room entry / exit before pre-anaesthesia entry & 309 / 2 of 14\,442 pairs. \\
\bottomrule
\end{tabularx}
\end{table}

Administrative spreadsheet completeness therefore does not imply valid
clinical timing. Zero-coded endpoints and inverted intervals require explicit
handling, not an automatic 24-hour correction. High-cardinality coding
(617 principal diagnoses, 570 principal procedures and 1433 intervention-type
labels) motivates category harmonisation. The audit's coded-sex categories
require confirmation before interpretation. The full aggregate source audit,
including variable families, population summaries and detailed LOS discrepancy
counts, is preserved in the pinned contribution archive. No patient or
professional identifiers are reproduced.
